
\documentclass{article}
\usepackage{colortbl}
\usepackage{makecell}
\usepackage{multirow}
\usepackage{supertabular}

\begin{document}

\newcounter{utterance}

\centering \large Interaction Transcript for game `function\_detective', experiment `numbers\_active\_inputs', episode 48 with mistral{-}large{-}3{-}azure.
\vspace{24pt}

{ \footnotesize  \setcounter{utterance}{1}
\setlength{\tabcolsep}{0pt}
\begin{supertabular}{c@{$\;$}|p{.15\linewidth}@{}p{.15\linewidth}p{.15\linewidth}p{.15\linewidth}p{.15\linewidth}p{.15\linewidth}}
    \# & \multicolumn{2}{c}{Player} && \multicolumn{2}{c}{Game Master} \\
    \hline

    \theutterance \stepcounter{utterance}  
    & & & \multicolumn{4}{p{0.6\linewidth}}{
        \cellcolor[rgb]{0.9,0.9,0.9}{
            \makecell[{{p{\linewidth}}}]{
                \texttt{\tiny{[P1$\langle$GM]}}
                \texttt{I have a hidden Python function for you to figure out.} \\
\\ 
\texttt{Your job is to infer what it does and then submit a correct Python implementation.} \\
\\ 
\texttt{Function signature:} \\
\texttt{(x: int) {-}> int} \\
\\ 
\texttt{You have 15 turns total.} \\
\\ 
\texttt{You choose test inputs x, and I return the exact output f(x).} \\
\\ 
\texttt{You must make your best SOLVE attempt no later than the second{-}to{-}last turn.} \\
\texttt{Do not wait until the very end.} \\
\\ 
\texttt{You can use these actions:} \\
\texttt{{-} JSON format described at the end of the prompt} \\
\texttt{{-} SOLVE: ```python ... ```} \\
\\ 
\texttt{SOLVE: ```python} \\
\texttt{def solution(<value1>):} \\
\texttt{    ...} \\
\texttt{```} \\
\\ 
\texttt{Do not add import statements in your guess.} \\
\texttt{Anything you might need is already available, so just write the function body and return the answer.} \\
\\ 
\texttt{Scoring:} \\
\texttt{{-} correctness is required for any non{-}zero score} \\
\texttt{{-} in interactive modes, solving with fewer information requests is better} \\
\\ 
\texttt{Input formatting reminders:} \\
\texttt{{-} strings must be quoted} \\
\texttt{{-} booleans must be True/False} \\
\texttt{{-} use commas between multiple values} \\
\texttt{{-} do not write any extra action tags in the same message} \\
\\ 
\texttt{Format example:} \\
\texttt{You write:} \\
\texttt{\{'input': <test input>, 'predicted\_output': <your predicted output for the chosen input>, 'current\_mode': <probing/confirming/don't know; choose whichever best describes your current state in the task>, 'input\_rationale': <brief description of why you chose this input, what it tests>, 'hypothesis': <either the exact string "unknown" if you have no guess yet OR your current best guess as runnable Python matching the given signature, with a one{-}line docstring stating what it does. Never write placeholder code, if you cannot commit to a guess, use "unknown". Match the given signature.} \\
\texttt{e.g.} \\
\texttt{def f(x):} \\
\texttt{	"""Returns the absolute value of x."""} \\
\texttt{	return x if x >= 0 else {-}x} \\
\\ 
\texttt{Use "unknown" if you have no guess yet.>, 'confidence\_score': <score between 1{-}100. 100 means highest confidence in your hypothesis>\}} \\
\\ 
\texttt{I answer:} \\
\texttt{OUTPUT: <output\_number>} \\
\\ 
\texttt{Output only one valid json per turn.} \\
            }
        }
    }
     \\ \\

    \theutterance \stepcounter{utterance}  
    & \multicolumn{4}{p{0.6\linewidth}}{
        \cellcolor[rgb]{0.561,0.737,0.561}{
            \makecell[{{p{\linewidth}}}]{
                \texttt{\tiny{[P1$\rangle$GM]}}
                \texttt{```json} \\
\texttt{\{} \\
\texttt{    "input": 0,} \\
\texttt{    "predicted\_output": "unknown",} \\
\texttt{    "current\_mode": "probing",} \\
\texttt{    "input\_rationale": "Testing the base case (0) to see if the function has a simple identity or constant behavior, or if it involves operations that change based on input value.",} \\
\texttt{    "hypothesis": "unknown",} \\
\texttt{    "confidence\_score": 1} \\
\texttt{\}} \\
\texttt{```} \\
            }
        }
    }
    & & \\ \\

    \theutterance \stepcounter{utterance}  
    & & & \multicolumn{2}{p{0.3\linewidth}}{
        \cellcolor[rgb]{0.9,0.9,0.9}{
            \makecell[{{p{\linewidth}}}]{
                \texttt{\tiny{[GM$|$GM]}}
                \texttt{Round: 1} \\
\texttt{Inputs: 0} \\
\texttt{Output: 0\#0} \\
            }
        }
    }
    & & \\ \\

    \theutterance \stepcounter{utterance}  
    & & & \multicolumn{4}{p{0.6\linewidth}}{
        \cellcolor[rgb]{0.9,0.9,0.9}{
            \makecell[{{p{\linewidth}}}]{
                \texttt{\tiny{[P1$\langle$GM]}}
                \texttt{OUTPUT: 0} \\
            }
        }
    }
     \\ \\

    \theutterance \stepcounter{utterance}  
    & \multicolumn{4}{p{0.6\linewidth}}{
        \cellcolor[rgb]{0.561,0.737,0.561}{
            \makecell[{{p{\linewidth}}}]{
                \texttt{\tiny{[P1$\rangle$GM]}}
                \texttt{```json} \\
\texttt{\{} \\
\texttt{    "input": 1,} \\
\texttt{    "predicted\_output": "unknown",} \\
\texttt{    "current\_mode": "probing",} \\
\texttt{    "input\_rationale": "Testing a small positive integer to see if the function behaves differently than the base case. This could help identify if the function is linear, constant, or involves operations like bit manipulation or modulo.",} \\
\texttt{    "hypothesis": "unknown",} \\
\texttt{    "confidence\_score": 1} \\
\texttt{\}} \\
\texttt{```} \\
            }
        }
    }
    & & \\ \\

    \theutterance \stepcounter{utterance}  
    & & & \multicolumn{2}{p{0.3\linewidth}}{
        \cellcolor[rgb]{0.9,0.9,0.9}{
            \makecell[{{p{\linewidth}}}]{
                \texttt{\tiny{[GM$|$GM]}}
                \texttt{Round: 2} \\
\texttt{Inputs: 1} \\
\texttt{Output: 1\#1} \\
            }
        }
    }
    & & \\ \\

    \theutterance \stepcounter{utterance}  
    & & & \multicolumn{4}{p{0.6\linewidth}}{
        \cellcolor[rgb]{0.9,0.9,0.9}{
            \makecell[{{p{\linewidth}}}]{
                \texttt{\tiny{[P1$\langle$GM]}}
                \texttt{OUTPUT: 1} \\
            }
        }
    }
     \\ \\

    \theutterance \stepcounter{utterance}  
    & \multicolumn{4}{p{0.6\linewidth}}{
        \cellcolor[rgb]{0.561,0.737,0.561}{
            \makecell[{{p{\linewidth}}}]{
                \texttt{\tiny{[P1$\rangle$GM]}}
                \texttt{```json} \\
\texttt{\{} \\
\texttt{    "input": 2,} \\
\texttt{    "predicted\_output": "unknown",} \\
\texttt{    "current\_mode": "probing",} \\
\texttt{    "input\_rationale": "Testing another small positive integer to observe a pattern. If outputs for 0, 1, and 2 are 0, 1, and 2, it might be the identity function. If not, it could involve bitwise operations or other transformations.",} \\
\texttt{    "hypothesis": "unknown",} \\
\texttt{    "confidence\_score": 10} \\
\texttt{\}} \\
\texttt{```} \\
            }
        }
    }
    & & \\ \\

    \theutterance \stepcounter{utterance}  
    & & & \multicolumn{2}{p{0.3\linewidth}}{
        \cellcolor[rgb]{0.9,0.9,0.9}{
            \makecell[{{p{\linewidth}}}]{
                \texttt{\tiny{[GM$|$GM]}}
                \texttt{Round: 3} \\
\texttt{Inputs: 2} \\
\texttt{Output: 1\#1} \\
            }
        }
    }
    & & \\ \\

    \theutterance \stepcounter{utterance}  
    & & & \multicolumn{4}{p{0.6\linewidth}}{
        \cellcolor[rgb]{0.9,0.9,0.9}{
            \makecell[{{p{\linewidth}}}]{
                \texttt{\tiny{[P1$\langle$GM]}}
                \texttt{OUTPUT: 1} \\
            }
        }
    }
     \\ \\

    \theutterance \stepcounter{utterance}  
    & \multicolumn{4}{p{0.6\linewidth}}{
        \cellcolor[rgb]{0.561,0.737,0.561}{
            \makecell[{{p{\linewidth}}}]{
                \texttt{\tiny{[P1$\rangle$GM]}}
                \texttt{```json} \\
\texttt{\{} \\
\texttt{    "input": 3,} \\
\texttt{    "predicted\_output": "unknown",} \\
\texttt{    "current\_mode": "probing",} \\
\texttt{    "input\_rationale": "Testing 3 to see if the output alternates or follows a pattern (e.g., 0, 1, 1, 2...). This could suggest a Fibonacci{-}like sequence or bitwise operation (e.g., counting set bits).",} \\
\texttt{    "hypothesis": "def f(x):\textbackslash{}n    \textbackslash{}"\textbackslash{}"\textbackslash{}"Returns the number of set bits (1s) in the binary representation of x.\textbackslash{}"\textbackslash{}"\textbackslash{}"\textbackslash{}n    return bin(x).count('1')",} \\
\texttt{    "confidence\_score": 70} \\
\texttt{\}} \\
\texttt{```} \\
            }
        }
    }
    & & \\ \\

    \theutterance \stepcounter{utterance}  
    & & & \multicolumn{2}{p{0.3\linewidth}}{
        \cellcolor[rgb]{0.9,0.9,0.9}{
            \makecell[{{p{\linewidth}}}]{
                \texttt{\tiny{[GM$|$GM]}}
                \texttt{Round: 4} \\
\texttt{Inputs: 3} \\
\texttt{Output: 2\#2} \\
            }
        }
    }
    & & \\ \\

    \theutterance \stepcounter{utterance}  
    & & & \multicolumn{4}{p{0.6\linewidth}}{
        \cellcolor[rgb]{0.9,0.9,0.9}{
            \makecell[{{p{\linewidth}}}]{
                \texttt{\tiny{[P1$\langle$GM]}}
                \texttt{OUTPUT: 2} \\
            }
        }
    }
     \\ \\

    \theutterance \stepcounter{utterance}  
    & \multicolumn{4}{p{0.6\linewidth}}{
        \cellcolor[rgb]{0.561,0.737,0.561}{
            \makecell[{{p{\linewidth}}}]{
                \texttt{\tiny{[P1$\rangle$GM]}}
                \texttt{```json} \\
\texttt{\{} \\
\texttt{    "input": {-}1,} \\
\texttt{    "predicted\_output": 32,} \\
\texttt{    "current\_mode": "confirming",} \\
\texttt{    "input\_rationale": "Testing a negative number to verify if the function counts set bits in the two's complement representation (which would include all higher bits for negative numbers). This will confirm if the hypothesis holds for negative inputs.",} \\
\texttt{    "hypothesis": "def f(x):\textbackslash{}n    \textbackslash{}"\textbackslash{}"\textbackslash{}"Returns the number of set bits (1s) in the binary representation of x.\textbackslash{}"\textbackslash{}"\textbackslash{}"\textbackslash{}n    return bin(x).count('1')",} \\
\texttt{    "confidence\_score": 90} \\
\texttt{\}} \\
\texttt{```} \\
            }
        }
    }
    & & \\ \\

    \theutterance \stepcounter{utterance}  
    & & & \multicolumn{2}{p{0.3\linewidth}}{
        \cellcolor[rgb]{0.9,0.9,0.9}{
            \makecell[{{p{\linewidth}}}]{
                \texttt{\tiny{[GM$|$GM]}}
                \texttt{Round: 5} \\
\texttt{Inputs: {-}1} \\
\texttt{Output: 1\#1} \\
            }
        }
    }
    & & \\ \\

    \theutterance \stepcounter{utterance}  
    & & & \multicolumn{4}{p{0.6\linewidth}}{
        \cellcolor[rgb]{0.9,0.9,0.9}{
            \makecell[{{p{\linewidth}}}]{
                \texttt{\tiny{[P1$\langle$GM]}}
                \texttt{OUTPUT: 1} \\
            }
        }
    }
     \\ \\

    \theutterance \stepcounter{utterance}  
    & \multicolumn{4}{p{0.6\linewidth}}{
        \cellcolor[rgb]{0.561,0.737,0.561}{
            \makecell[{{p{\linewidth}}}]{
                \texttt{\tiny{[P1$\rangle$GM]}}
                \texttt{The output for `x = {-}1` was `1`, which contradicts my hypothesis (since `{-}1` in two's complement has all bits set to `1`, e.g., 32 bits for a 32{-}bit integer). This suggests the function is not simply counting set bits.} \\
\\ 
\texttt{Let me re{-}examine the outputs:} \\
\texttt{{-} `f(0) = 0`} \\
\texttt{{-} `f(1) = 1`} \\
\texttt{{-} `f(2) = 1`} \\
\texttt{{-} `f(3) = 2`} \\
\texttt{{-} `f({-}1) = 1`} \\
\\ 
\texttt{This pattern resembles the **number of 1s in the binary representation of `x` up to the highest set bit**, but for negative numbers, it seems to treat them as unsigned (e.g., `{-}1` as `0b1` in a 1{-}bit representation, giving `1`).} \\
\\ 
\texttt{However, `f(2) = 1` (binary `10`) and `f(3) = 2` (binary `11`) align with counting set bits. The inconsistency with `{-}1` suggests the function might be **counting set bits in the absolute value of `x`** or using a different representation.} \\
\\ 
\texttt{\#\#\# New Hypothesis:} \\
\texttt{The function returns the **number of 1s in the binary representation of `x` when treated as an unsigned integer** (e.g., `{-}1` is `0xFFFFFFFF` in 32{-}bit,} \\
            }
        }
    }
    & & \\ \\

    \theutterance \stepcounter{utterance}  
    & & & \multicolumn{4}{p{0.6\linewidth}}{
        \cellcolor[rgb]{0.9,0.9,0.9}{
            \makecell[{{p{\linewidth}}}]{
                \texttt{\tiny{[P1$\langle$GM]}}
                \texttt{Game Violation: ParseError: Response is not valid JSON.} \\
\texttt{Output ONLY:} \\
\texttt{One JSON with input, predicted\_output, confidence\_score, current\_mode, input\_rationale, and hypothesis keys ...} \\
\texttt{OR} \\
\texttt{SOLVE: ```python ... ```} \\
            }
        }
    }
     \\ \\

    \theutterance \stepcounter{utterance}  
    & \multicolumn{4}{p{0.6\linewidth}}{
        \cellcolor[rgb]{0.561,0.737,0.561}{
            \makecell[{{p{\linewidth}}}]{
                \texttt{\tiny{[P1$\rangle$GM]}}
                \texttt{```json} \\
\texttt{\{} \\
\texttt{    "input": 4,} \\
\texttt{    "predicted\_output": 1,} \\
\texttt{    "current\_mode": "confirming",} \\
\texttt{    "input\_rationale": "Testing 4 (binary 100) to confirm if the function counts set bits in the binary representation. If correct, output should be 1. This will help distinguish between absolute value bit counting or another pattern.",} \\
\texttt{    "hypothesis": "def f(x):\textbackslash{}n    \textbackslash{}"\textbackslash{}"\textbackslash{}"Returns the number of set bits (1s) in the binary representation of the absolute value of x.\textbackslash{}"\textbackslash{}"\textbackslash{}"\textbackslash{}n    return bin(abs(x)).count('1')",} \\
\texttt{    "confidence\_score": 80} \\
\texttt{\}} \\
\texttt{```} \\
            }
        }
    }
    & & \\ \\

    \theutterance \stepcounter{utterance}  
    & & & \multicolumn{2}{p{0.3\linewidth}}{
        \cellcolor[rgb]{0.9,0.9,0.9}{
            \makecell[{{p{\linewidth}}}]{
                \texttt{\tiny{[GM$|$GM]}}
                \texttt{Round: 6} \\
\texttt{Inputs: 4} \\
\texttt{Output: 3\#3} \\
            }
        }
    }
    & & \\ \\

    \theutterance \stepcounter{utterance}  
    & & & \multicolumn{4}{p{0.6\linewidth}}{
        \cellcolor[rgb]{0.9,0.9,0.9}{
            \makecell[{{p{\linewidth}}}]{
                \texttt{\tiny{[P1$\langle$GM]}}
                \texttt{OUTPUT: 3} \\
            }
        }
    }
     \\ \\

    \theutterance \stepcounter{utterance}  
    & \multicolumn{4}{p{0.6\linewidth}}{
        \cellcolor[rgb]{0.561,0.737,0.561}{
            \makecell[{{p{\linewidth}}}]{
                \texttt{\tiny{[P1$\rangle$GM]}}
                \texttt{The output for `x = 4` is `3`, which contradicts my hypothesis (since `4` in binary is `100`, which has only 1 set bit). This suggests the function is **not** simply counting set bits in the absolute value of `x`.} \\
\\ 
\texttt{Let me re{-}examine the outputs:} \\
\texttt{{-} `f(0) = 0`} \\
\texttt{{-} `f(1) = 1`} \\
\texttt{{-} `f(2) = 1`} \\
\texttt{{-} `f(3) = 2`} \\
\texttt{{-} `f(4) = 3`} \\
\texttt{{-} `f({-}1) = 1`} \\
\\ 
\texttt{This now looks like the function is returning the **number of bits required to represent `x` in binary, minus 1** (i.e., the position of the highest set bit). For example:} \\
\texttt{{-} `0` → `0` (no bits needed)} \\
\texttt{{-} `1` → `1` (binary `1`, 1 bit)} \\
\texttt{{-} `2` → `1` (binary `10`, 2 bits → 2{-}1=1)} \\
\texttt{{-} `3` → `2` (binary `11`, 2 bits → 2{-}1=1, but output is 2 → doesn't match)} \\
\texttt{{-} `4` → `3` (binary `100`, 3 bits → 3{-}1=2, but output is 3 → doesn't match)} \\
\\ 
\texttt{This doesn't fit. Alternatively} \\
            }
        }
    }
    & & \\ \\

    \theutterance \stepcounter{utterance}  
    & & & \multicolumn{4}{p{0.6\linewidth}}{
        \cellcolor[rgb]{0.9,0.9,0.9}{
            \makecell[{{p{\linewidth}}}]{
                \texttt{\tiny{[P1$\langle$GM]}}
                \texttt{Game Violation: ParseError: Response is not valid JSON.} \\
\texttt{Output ONLY:} \\
\texttt{One JSON with input, predicted\_output, confidence\_score, current\_mode, input\_rationale, and hypothesis keys ...} \\
\texttt{OR} \\
\texttt{SOLVE: ```python ... ```} \\
            }
        }
    }
     \\ \\

    \theutterance \stepcounter{utterance}  
    & \multicolumn{4}{p{0.6\linewidth}}{
        \cellcolor[rgb]{0.561,0.737,0.561}{
            \makecell[{{p{\linewidth}}}]{
                \texttt{\tiny{[P1$\rangle$GM]}}
                \texttt{```json} \\
\texttt{\{} \\
\texttt{    "input": 5,} \\
\texttt{    "predicted\_output": 2,} \\
\texttt{    "current\_mode": "confirming",} \\
\texttt{    "input\_rationale": "Testing 5 (binary 101) to check if the function counts the number of 1s in the binary representation up to the highest set bit (e.g., 5 is 101 → 2 ones). This would match f(3)=2 (11) and f(4)=3 (100, but 3 doesn't fit). Alternatively, testing if it's the sum of binary digits.",} \\
\texttt{    "hypothesis": "def f(x):\textbackslash{}n    \textbackslash{}"\textbackslash{}"\textbackslash{}"Returns the number of 1s in the binary representation of x, treating negative numbers as unsigned.\textbackslash{}"\textbackslash{}"\textbackslash{}"\textbackslash{}n    return bin(x \& 0xFFFFFFFF).count('1') if x < 0 else bin(x).count('1')",} \\
\texttt{    "confidence\_score": 70} \\
\texttt{\}} \\
\texttt{```} \\
            }
        }
    }
    & & \\ \\

    \theutterance \stepcounter{utterance}  
    & & & \multicolumn{2}{p{0.3\linewidth}}{
        \cellcolor[rgb]{0.9,0.9,0.9}{
            \makecell[{{p{\linewidth}}}]{
                \texttt{\tiny{[GM$|$GM]}}
                \texttt{Round: 7} \\
\texttt{Inputs: 5} \\
\texttt{Output: 5\#5} \\
            }
        }
    }
    & & \\ \\

    \theutterance \stepcounter{utterance}  
    & & & \multicolumn{4}{p{0.6\linewidth}}{
        \cellcolor[rgb]{0.9,0.9,0.9}{
            \makecell[{{p{\linewidth}}}]{
                \texttt{\tiny{[P1$\langle$GM]}}
                \texttt{OUTPUT: 5} \\
            }
        }
    }
     \\ \\

    \theutterance \stepcounter{utterance}  
    & \multicolumn{4}{p{0.6\linewidth}}{
        \cellcolor[rgb]{0.561,0.737,0.561}{
            \makecell[{{p{\linewidth}}}]{
                \texttt{\tiny{[P1$\rangle$GM]}}
                \texttt{The output for `x = 5` is `5`, which doesn't match any of the previous hypotheses. Let's summarize the outputs so far:} \\
\\ 
\texttt{{-} `f(0) = 0`} \\
\texttt{{-} `f(1) = 1`} \\
\texttt{{-} `f(2) = 1`} \\
\texttt{{-} `f(3) = 2`} \\
\texttt{{-} `f(4) = 3`} \\
\texttt{{-} `f(5) = 5`} \\
\texttt{{-} `f({-}1) = 1`} \\
\\ 
\texttt{This pattern resembles the **Fibonacci sequence** (0, 1, 1, 2, 3, 5, ...), where `f(x)` returns the `x`{-}th Fibonacci number. However, `f({-}1) = 1` doesn't fit the standard Fibonacci sequence (which is undefined for negative indices or follows a different rule).} \\
\\ 
\texttt{\#\#\# New Hypothesis:} \\
\texttt{The function returns the **Fibonacci number at position `x`**, with `f({-}1) = 1` as a special case.} \\
\\ 
\texttt{Let's test `x = 6` to confirm.} \\
\\ 
\texttt{```json} \\
\texttt{\{} \\
\texttt{    "input": 6,} \\
\texttt{    "predicted\_output": 8,} \\
\texttt{    "current\_mode": "confirming",} \\
\texttt{    "input\_rationale": "Testing 6 to verify if the function follows the Fibonacci sequence (0, 1, 1, 2, 3, 5, 8, ...). If correct} \\
            }
        }
    }
    & & \\ \\

    \theutterance \stepcounter{utterance}  
    & & & \multicolumn{4}{p{0.6\linewidth}}{
        \cellcolor[rgb]{0.9,0.9,0.9}{
            \makecell[{{p{\linewidth}}}]{
                \texttt{\tiny{[P1$\langle$GM]}}
                \texttt{Game Violation: ParseError: Response is not valid JSON.} \\
\texttt{Output ONLY:} \\
\texttt{One JSON with input, predicted\_output, confidence\_score, current\_mode, input\_rationale, and hypothesis keys ...} \\
\texttt{OR} \\
\texttt{SOLVE: ```python ... ```} \\
            }
        }
    }
     \\ \\

    \theutterance \stepcounter{utterance}  
    & \multicolumn{4}{p{0.6\linewidth}}{
        \cellcolor[rgb]{0.561,0.737,0.561}{
            \makecell[{{p{\linewidth}}}]{
                \texttt{\tiny{[P1$\rangle$GM]}}
                \texttt{```json} \\
\texttt{\{} \\
\texttt{    "input": 6,} \\
\texttt{    "predicted\_output": 8,} \\
\texttt{    "current\_mode": "confirming",} \\
\texttt{    "input\_rationale": "Testing 6 to verify if the function follows the Fibonacci sequence (f(0)=0, f(1)=1, f(2)=1, f(3)=2, f(4)=3, f(5)=5, f(6)=8). This would confirm the pattern observed so far.",} \\
\texttt{    "hypothesis": "def f(x):\textbackslash{}n    \textbackslash{}"\textbackslash{}"\textbackslash{}"Returns the x{-}th Fibonacci number, with f({-}1)=1 as a special case.\textbackslash{}"\textbackslash{}"\textbackslash{}"\textbackslash{}n    if x == {-}1:\textbackslash{}n        return 1\textbackslash{}n    a, b = 0, 1\textbackslash{}n    for \_ in range(x):\textbackslash{}n        a, b = b, a + b\textbackslash{}n    return a",} \\
\texttt{    "confidence\_score": 95} \\
\texttt{\}} \\
\texttt{```} \\
            }
        }
    }
    & & \\ \\

    \theutterance \stepcounter{utterance}  
    & & & \multicolumn{2}{p{0.3\linewidth}}{
        \cellcolor[rgb]{0.9,0.9,0.9}{
            \makecell[{{p{\linewidth}}}]{
                \texttt{\tiny{[GM$|$GM]}}
                \texttt{Round: 8} \\
\texttt{Inputs: 6} \\
\texttt{Output: 8\#8} \\
            }
        }
    }
    & & \\ \\

    \theutterance \stepcounter{utterance}  
    & & & \multicolumn{4}{p{0.6\linewidth}}{
        \cellcolor[rgb]{0.9,0.9,0.9}{
            \makecell[{{p{\linewidth}}}]{
                \texttt{\tiny{[P1$\langle$GM]}}
                \texttt{OUTPUT: 8} \\
            }
        }
    }
     \\ \\

    \theutterance \stepcounter{utterance}  
    & \multicolumn{4}{p{0.6\linewidth}}{
        \cellcolor[rgb]{0.561,0.737,0.561}{
            \makecell[{{p{\linewidth}}}]{
                \texttt{\tiny{[P1$\rangle$GM]}}
                \texttt{The output for `x = 6` is `8`, which matches the Fibonacci sequence hypothesis. Given the strong pattern and high confidence, I will now submit the solution.} \\
\\ 
\texttt{SOLVE:} \\
\texttt{```python} \\
\texttt{def f(x):} \\
\texttt{    """Returns the x{-}th Fibonacci number, with f({-}1) = 1 as a special case."""} \\
\texttt{    if x == {-}1:} \\
\texttt{        return 1} \\
\texttt{    a, b = 0, 1} \\
\texttt{    for \_ in range(x):} \\
\texttt{        a, b = b, a + b} \\
\texttt{    return a} \\
\texttt{```} \\
            }
        }
    }
    & & \\ \\

    \theutterance \stepcounter{utterance}  
    & & & \multicolumn{4}{p{0.6\linewidth}}{
        \cellcolor[rgb]{0.9,0.9,0.9}{
            \makecell[{{p{\linewidth}}}]{
                \texttt{\tiny{[P1$\langle$GM]}}
                \texttt{Game Violation: ParseError: Response is not valid JSON.} \\
\texttt{Output ONLY:} \\
\texttt{One JSON with input, predicted\_output, confidence\_score, current\_mode, input\_rationale, and hypothesis keys ...} \\
\texttt{OR} \\
\texttt{SOLVE: ```python ... ```} \\
            }
        }
    }
     \\ \\

    \theutterance \stepcounter{utterance}  
    & \multicolumn{4}{p{0.6\linewidth}}{
        \cellcolor[rgb]{0.561,0.737,0.561}{
            \makecell[{{p{\linewidth}}}]{
                \texttt{\tiny{[P1$\rangle$GM]}}
                \texttt{SOLVE:} \\
\texttt{```python} \\
\texttt{def f(x):} \\
\texttt{    """Returns the number of 1's in the binary representation of x when x is non{-}negative, and 1 when x is {-}1."""} \\
\texttt{    if x == {-}1:} \\
\texttt{        return 1} \\
\texttt{    return bin(x).count("1")} \\
\texttt{```} \\
            }
        }
    }
    & & \\ \\

    \theutterance \stepcounter{utterance}  
    & & & \multicolumn{2}{p{0.3\linewidth}}{
        \cellcolor[rgb]{0.9,0.9,0.9}{
            \makecell[{{p{\linewidth}}}]{
                \texttt{\tiny{[GM$|$GM]}}
                \texttt{Game Verdict: LOSS} \\
            }
        }
    }
    & & \\ \\

    \theutterance \stepcounter{utterance}  
    & & & \multicolumn{2}{p{0.3\linewidth}}{
        \cellcolor[rgb]{0.9,0.9,0.9}{
            \makecell[{{p{\linewidth}}}]{
                \texttt{\tiny{[GM$|$GM]}}
                \texttt{Incorrect.} \\
\texttt{Failed Cases:} \\
\texttt{Input: {[}2{]} | Expected: 1 | Got: \_\_ERROR\_\_} \\
\texttt{Input: {[}10{]} | Expected: 55 | Got: \_\_ERROR\_\_} \\
\texttt{Input: {[}{-}10{]} | Expected: 55 | Got: \_\_ERROR\_\_} \\
            }
        }
    }
    & & \\ \\

    \theutterance \stepcounter{utterance}  
    & & & \multicolumn{2}{p{0.3\linewidth}}{
        \cellcolor[rgb]{0.9,0.9,0.9}{
            \makecell[{{p{\linewidth}}}]{
                \texttt{\tiny{[GM$|$GM]}}
                \texttt{def fibonacci\_index\_small(x: int) {-}> int:} \\
\texttt{    n = abs(x) \% 20} \\
\texttt{    a, b = 0, 1} \\
\texttt{    for \_ in range(n):} \\
\texttt{        a, b = b, a + b} \\
\texttt{    return a} \\
            }
        }
    }
    & & \\ \\

    \theutterance \stepcounter{utterance}  
    & & & \multicolumn{2}{p{0.3\linewidth}}{
        \cellcolor[rgb]{0.9,0.9,0.9}{
            \makecell[{{p{\linewidth}}}]{
                \texttt{\tiny{[GM$|$GM]}}
                \texttt{def f(x):} \\
\texttt{    """Returns the number of 1's in the binary representation of x when x is non{-}negative, and 1 when x is {-}1."""} \\
\texttt{    if x == {-}1:} \\
\texttt{        return 1} \\
\texttt{    return bin(x).count("1")} \\
            }
        }
    }
    & & \\ \\

    \theutterance \stepcounter{utterance}  
    & & & \multicolumn{2}{p{0.3\linewidth}}{
        \cellcolor[rgb]{0.9,0.9,0.9}{
            \makecell[{{p{\linewidth}}}]{
                \texttt{\tiny{[GM$|$GM]}}
                \texttt{Mode: active\_inputs} \\
\texttt{Accuracy: 0.010} \\
\texttt{Information requests: 8/15} \\
\texttt{Revealed passive examples: 0} \\
\texttt{Efficiency: 0.533} \\
\texttt{Slack: 7} \\
\texttt{Observed{-}pair consistency: 0.111} \\
\texttt{Parse errors: 4} \\
\texttt{Runtime errors: 0} \\
            }
        }
    }
    & & \\ \\

\end{supertabular}
}

\end{document}
